% R. Brown's translations of two letters to L. Breen, with his note and comments (cote 134-2): archive cote 134-2, site page 117 (PDF page 118).
% comments. Draft compared: none.
% AI-transcribed by Claude Opus 5.5 (claude-opus-5-5, adaptive thinking, effort high), Message Batches API, batch brown_full (msgbatch_01Md3i5yK1PxkpJn5GHfEExE), 2026-11-13; not yet checked by eye.
\scanpage{134-2/117}
\uline{COMMENTS}

1. The topic of duality has progressed somewhat since 1975, see for
example

L. Begueri \hand{Memoir SMF}
\margin{demander réf. : Breen.}

M. Bester \hand{Thesis, U of Michigan}

(These references are due to L. Breen.)

2. The work of Mme Sinh referred to on p. 2 is as follows:

H.X. Sinh, \textit{Gr-catégories}, Thèse Doctorat d'\over{E}{É}tat Université Paris VII
(1975).

See also

N.\deleted{\unsure{S}} \inserted{Saavedra} Rivano, \textit{Catégories tanakiennes}, Springer Lecture
Notes in Mathematics 265 (1972).

Related to these notions is that of +-groupoid (=coherent
symmetric monoidal groupoid with strict associativity and strict
zero object) whose theory is developed also in:

M. Takeuchi, 'On Villamayor and Zelinsky's long exact sequence',
Mem. Amer. Math. Soc. 33 (1981) No 249.

3. The description of groups in (1-cat) given on p. 2-3 is in effect an
equivalence between G-groupoids (= group objects in the category of
groupoids) and crossed modules. This is described also in:

R. Brown and C.B. Spencer, 'G-groupoids, crossed modules and the
fundamental groupoid of a topological group', Proc. Kon. Akad.
v. Weten. 79(1976), 296 - 302.
and is believed to have been known to Verdier in 1965.

4. The equivalence mentioned in the last comment has been generalised to

\mspage{118}
